% Section 9: Contraction in four guises} | Part: 1 | Depends on: see docs/REWRITE_PLAN.md | Status: todo | Original lines: 1668-1722
%!TEX root = ../../../main.tex
\subsection{Contraction in four guises}\label{sec:contr4}

The word ``contraction'' has been used four times; they are one idea at four levels of structure.

\begin{center}
\small
\renewcommand{\arraystretch}{1.35}
\begin{tabular}{@{}p{3.4cm}p{3.5cm}p{4.2cm}p{2.2cm}@{}}
\toprule
\textbf{Guise} & \textbf{Acts on} & \textbf{Formula} & \textbf{Metric?}\\
\midrule
Evaluation $\langle\omega,v\rangle$ & $V^*\times V\to\R$ & $\omega_iv^i$ & no\\
Tensor contraction & $T^{\ell}_{k}\to T^{\ell-1}_{k-1}$ & trace $A^i{}_i$; $f_iv^i$; Ricci $R^i{}_{ijk}$ & no\\
Interior product $\iota_v$ & $k$-forms $\to(k-1)$-forms & $(\iota_v\alpha)(w_\bullet)=\alpha(v,w_\bullet)$ & no\\
Left contraction $a\lc A$ & $k$-vectors $\to(k-1)$-vectors & $a\lc(b\wedge c)=(a\!\cdot\! b)c-(a\!\cdot\! c)b$ & yes ($\cdot$ is $g$)\\
\bottomrule
\end{tabular}
\end{center}

The first three need only $V$ and $V^*$. The fourth acts on vectors and needs a metric, but it is the \emph{same} operation as the third once a
metric is available to translate between $V$ and $V^*$:

\begin{proposition}[left contraction $=$ interior product under $\flat$]\label{prop:lcisiota}
Extend $\flat:V\to V^*$ to an algebra homomorphism of exterior algebras, $(v_1\wedge\cdots\wedge v_k)^\flat=v_1^\flat\wedge\cdots\wedge v_k^\flat$. Then for every $a\in V$
and every multivector $A$,
\[
(a\lc A)^\flat=\iota_a\bigl(A^\flat\bigr).
\]
\end{proposition}
\begin{proof}
Both sides are antiderivations in $A$ of degree $-1$ (for $a\lc$, see the properties in Section~\ref{sec:ops}; for $\iota_a$, \eqref{eq:antider}). On a vector $b$ they agree:
$(a\lc b)^\flat=(a\cdot b)=g(a,b)=b^\flat(a)=\iota_a b^\flat$. Since multivectors are sums of wedges of vectors, an induction on the grade using the antiderivation rule shows that they agree everywhere.
\end{proof}

\begin{example}[the three-vector check]\label{ex:lcvsiota}
Take $a,b,c\in\R^3$ with the Euclidean metric. Then $(a\lc(b\wedge c))^\flat=\bigl((a\!\cdot\!b)c-(a\!\cdot\!c)b\bigr)^\flat=(a\!\cdot\!b)c^\flat-(a\!\cdot\!c)b^\flat$, and
$\iota_a(b^\flat\wedge c^\flat)=b^\flat(a)\,c^\flat-c^\flat(a)\,b^\flat$ by \eqref{eq:wedge11}; the two agree because $b^\flat(a)=a\cdot b$. In coordinates this is exactly
the planar computation of Example~\ref{ex:interiorplane}: $\iota_v(\dd x\wedge\dd y)=v^x\dd y-v^y\dd x$.
\end{example}

\begin{synthesis}[the same identities, three notations]
\renewcommand{\arraystretch}{1.3}
\begin{center}\small
\begin{tabular}{@{}p{3.6cm}p{4.0cm}p{4.0cm}@{}}
\toprule
\textbf{Fact} & \textbf{Forms} & \textbf{Multivectors}\\
\midrule
antiderivation & $\iota_v(\alpha\wedge\beta)=\iota_v\alpha\wedge\beta+(-1)^k\alpha\wedge\iota_v\beta$ & $a\lc(A_k\wedge B)=(a\lc A_k)\wedge B+(-1)^kA_k\wedge(a\lc B)$\\
nilpotent & $\iota_v\iota_v=0$ & $a\lc(a\lc A)=(a\wedge a)\lc A=0$\\
wedge is not metric & $\alpha\wedge\beta$ & $A\wedge B$ (grade part of $AB$, but defined without $g$)\\
\bottomrule
\end{tabular}
\end{center}
\end{synthesis}
